\subsection{The structure of finite fields}

PROPOSITION 1. --- Let K be a finite field with q elements.

\begin{enumerate}
\def\labelenumi{\alph{enumi})}
\item
  The characteristic of \(\mathbf{K}\) is a prime number \(p\) , and there exists an integer \(f \geqslant 1\) such that \(q = p^f\) .
\item
  The additive group of \(\mathbf{K}\) is a direct sum of \(f\) cyclic groups of order \(p\) .
\item
  The multiplicative group of \(\mathbf{K}\) is cyclic of order \(q - 1\) .
\end{enumerate}

Since Z is infinite and K finite, the unique ring homomorphism \(\varphi:Z\to K\) is not injective and its kernel is a prime ideal of Z not equal to 0. Therefore the characteristic of K is a prime number and K is an algebra over the field \(F_{p}\) of p elements (V, p.~2). Let f be the degree of K over \(F_{p}\) . If f were infinite, K would contain, for each integer \(n\geqslant0\) a subspace of dimension n over \(F_{p}\) , whence \(q\geqslant p^{n}\) , which is absurd. Therefore f is finite. The additive group of K is thus isomorphic to \((\mathbf{F}_{p})^{f}\) , whence the assertions a) and b).

Assertion c) follows by Lemma 1 (V, p.~78).

PROPOSITION 2. --- Let K be a finite field of q elements. The field K is a splitting field of the polynomial \(X^{q}-X\) of \(F_{p}[X]\) and it is the set of all roots of this polynomial.

For every \(x \neq 0\) in \(\mathbf{K}\) we have \(x^{q-1} = 1\) because \(\mathbf{K}^*\) is a finite group of order \(q - 1\) (I, p.~52). It follows that \(x^q = x\) for each \(x\) in \(\mathbf{K}\) . The polynomial \(\mathbf{X}^q - \mathbf{X}\) of \(\mathbf{F}_p[\mathbf{X}]\) is of degree \(q\) and it has \(q\) roots in \(\mathbf{K}\) , whence

\[
\mathbf {X} ^ {q} - \mathbf {X} = \prod_ {\xi \in K} (\mathbf {X} - \xi).\tag{1}
\]

Proposition 2 follows at once from this.

COROLLARY. --- Two finite fields of the same cardinal are isomorphic.

Let \(\mathbf{K}'\) be a finite field of \(q\) elements; its characteristic is a prime number \(p'\) dividing \(q = p^f\) , whence \(p' = p\) . Therefore \(\mathbf{K}'\) is a splitting field of the polynomial \(\mathbf{X}^q - \mathbf{X}\) in \(\mathbf{F}_p[\mathbf{X}]\) (Prop. 2), and so \(\mathbf{K}\) and \(\mathbf{K}'\) are isomorphic (V, p.~22, Cor.).

When \(\mathbf{K} = \mathbf{F}_p\) , Formula (1) reduces to the relation

\[
\mathbf {X} ^ {p} - \mathbf {X} \equiv \prod_ {i = 0} ^ {p - 1} (\mathbf {X} - i) \bmod p \mathbf {Z} [ \mathbf {X} ]\tag{2}
\]

in the polynomial ring Z{[}X{]}.

Formula (2) may also be written

\[
\mathbf {X} ^ {p - 1} - 1 \equiv \prod_ {i = 1} ^ {p - 1} (\mathbf {X} - i) \bmod p \mathbf {Z} [ \mathbf {X} ].\tag{3}
\]

In particular for X = 0 we obtain (« Wilson's formula »)

\[
(p - 1)! \equiv - 1 \bmod p.\tag{4}
\]

